\documentclass[11pt,a4paper]{article}
% main (standalone preamble for editing)
% vim: nowrap: spell:
% ══════════════════════════════════════════════════════════════════════════════
% Speed Maths --- shared preamble
% Included by every sheet and answer file via:  % main (standalone preamble for editing)
% vim: nowrap: spell:
% ══════════════════════════════════════════════════════════════════════════════
% Speed Maths --- shared preamble
% Included by every sheet and answer file via:  % main (standalone preamble for editing)
% vim: nowrap: spell:
% ══════════════════════════════════════════════════════════════════════════════
% Speed Maths --- shared preamble
% Included by every sheet and answer file via:  \input{../../shared/preamble}
% Editing THIS file changes the look of every sheet/answer across every pillar.
% ══════════════════════════════════════════════════════════════════════════════

\usepackage[a4paper, top=25mm, bottom=25mm, left=22mm, right=22mm]{geometry}
\usepackage{amsmath, amssymb}
\usepackage{enumitem}
\usepackage{booktabs}
\usepackage{array}
\usepackage{parskip}
\usepackage{microtype}
\usepackage{fancyhdr}
\usepackage{titlesec}
\usepackage{mdframed}
\usepackage{xcolor}
\usepackage[hidelinks]{hyperref}
\usepackage{graphicx}
\usepackage{tikz}
\usepackage{pgfplots}
\pgfplotsset{compat=1.18}

% ── Colours ───────────────────────────────────────────────────────────────────
\definecolor{darkgray}{gray}{0.15}
\definecolor{medgray}{gray}{0.45}
\definecolor{lightgray}{gray}{0.93}
\definecolor{boxgray}{gray}{0.88}

% ── Section formatting ──────────────────────────────────────────────────────────
\titleformat{\section}[block]
  {\normalfont\large\bfseries}{}
  {0pt}{}[\vspace{2pt}\hrule\vspace{6pt}]
\titlespacing*{\section}{0pt}{14pt}{4pt}

% ── List formatting ─────────────────────────────────────────────────────────────
\setlist[enumerate,1]{label=\textbf{\arabic*.}, leftmargin=2em, itemsep=10pt}

% ── Branding constants (edit here, not per-file) ────────────────────────────────
\newcommand{\SpeedExamLine}{\textbf{Speed Maths:} Competition \& Exam Prep}
\newcommand{\SpeedCredit}{Series by \href{https://www.linkedin.com/in/cerealdev/}{David Akinyele-Aje}}

% ── Header / footer ──────────────────────────────────────────────────────────────
% Usage (once, after \input{../../shared/preamble} and before \begin{document}):
%   \SpeedHeader{Algebra}{3}
% First arg  = pillar name shown as "Speed <Pillar> --- Daily Drill"
% Second arg = worksheet number
\pagestyle{fancy}
\newcommand{\SpeedHeader}[2]{%
  \fancyhf{}%
  \renewcommand{\headrulewidth}{0.4pt}%
  \setlength{\headheight}{14pt}%
  \fancyhead[L]{\small\textbf{Speed #1 --- Daily Drill}}%
  \fancyhead[R]{\small Worksheet \##2}%
  \fancyfoot[L]{\small \SpeedExamLine}%
  \fancyfoot[C]{\small\thepage}%
  \fancyfoot[R]{\small \SpeedCredit}%
  \renewcommand{\footrulewidth}{0.4pt}%
}

% ── Title block (used at the top of every sheet AND answer file) ───────────────
% Usage: \SpeedTitleBlock{Daily Algebra Drill \#3}{Jane Doe}
% %% AGENTS: When generating a new sheet, ask the human contributor for the name they want credited in argument 2.
% For answer files, pass the "--- Answers \& Investigations" suffix in the arg.
\newcommand{\SpeedTitleBlock}[2]{%
  \begin{center}
    {\LARGE\bfseries #1}\\[4pt]
    {\normalsize\itshape Sheet by #2}\\[6pt]
    \hrule\vspace{4pt}
    \begin{tabular}{llcc}
      \toprule
      \textbf{Section} & \textbf{Topic} & \textbf{Questions} & \textbf{Target time}\\
      \midrule
      A & Rapid Recognition          & 10 & 2 min 30 sec \\
      B & Manipulation Drills        & 10 & 8 min        \\
      C & Substitution \& Structure  & 8  & 10 min       \\
      D & Challenge                  & 5  & 15 min       \\
      \bottomrule
    \end{tabular}
    \vspace{4pt}
    \hrule
  \end{center}
}

% ── Sheet metadata: topic + the day's new tools ────────────────────────────────
% Usage (immediately after \SpeedTitleBlock, sheets only):
%   \SpeedMeta{Expanding, factorising, and the identity toolkit}
%             {difference of squares, sum/product form, completing the square}
%
% Arg 1 = the sheet's topic in one line. The website lists this instead of
%         "Sheet 03", so write the thing a reader would actually search for.
% Arg 2 = comma-separated, at most five: the tools this sheet introduces that
%         earlier sheets in the pillar did not. These become the website's tags.
%         Leave out anything the reader already met — the tags are meant to say
%         what is new today, not to summarise the sheet.
%
% Both are parsed straight out of this file by tools/build_website.py, so the
% sheet stays the single source of truth and the site cannot drift from it.
%
% This renders NOTHING. It is a declaration for the build script, not a piece of
% the sheet's layout: adding it to a sheet must not change that sheet's PDF, so
% the 25 existing sheets can be annotated without a single one of them being
% reissued. If the toolkit should also appear on the page, that is a separate
% decision about the sheet design — make it deliberately, here, not by leaving
% this macro printing something nobody asked it to.
\newcommand{\SpeedMeta}[2]{}

% ── Answer / Method / Investigate-further boxes (answer files only) ────────────
\newmdenv[
  backgroundcolor=lightgray,
  linecolor=medgray,
  linewidth=0.5pt,
  leftmargin=0pt, rightmargin=0pt,
  innerleftmargin=8pt, innerrightmargin=8pt,
  innertopmargin=5pt, innerbottommargin=5pt,
  skipabove=4pt, skipbelow=4pt
]{ansbox}

\newmdenv[
  backgroundcolor=white,
  linecolor=darkgray,
  linewidth=0.8pt,
  leftline=true, rightline=false, topline=false, bottomline=false,
  leftmargin=0pt, rightmargin=0pt,
  innerleftmargin=8pt, innerrightmargin=6pt,
  innertopmargin=4pt, innerbottommargin=4pt,
  skipabove=4pt, skipbelow=6pt
]{invbox}

\newcommand{\ans}[1]{%
  \begin{ansbox}
    \textbf{Answer:} #1
  \end{ansbox}}

\newcommand{\method}[1]{%
  {\small\textbf{Method:} #1}\par}

\newcommand{\inv}[1]{%
  \begin{invbox}
    \textbf{\small Investigate further:} {\small #1}
  \end{invbox}}

% ── Misc helpers ────────────────────────────────────────────────────────────────
\newcommand{\qn}[1]{\textbf{#1.}\;}

% Mistake-linking: attach a Top 5 Patterns / Common Traps bullet to the question
% label(s) in this sheet where it actually shows up, e.g. \seealso{B6, D2}
\newcommand{\seealso}[1]{\hfill\textit{\footnotesize(see #1)}}

% ── Graph options macro ─────────────────────────────────────────────────────────
% Usage: \GraphOptions{tikz code for A}{tikz code for B}{tikz code for C}{tikz code for D}
\newcommand{\GraphOptions}[4]{%
  \begin{center}
    \begin{tabular}{cc}
      \textbf{A)} & \textbf{B)} \\
      #1 & #2 \\[10pt]
      \textbf{C)} & \textbf{D)} \\
      #3 & #4
    \end{tabular}
  \end{center}
}

% ── Closing motivational rule + cross-promo footer (sheets only) ────────────────
% Usage: \SpeedClosing{"Quote text."}
\newcommand{\SpeedClosing}[1]{%
  \vspace{20pt}
  \begin{center}
  \rule{0.6\textwidth}{0.4pt}\\[4pt]
  {\small\itshape #1}\\[4pt]
  \rule{0.6\textwidth}{0.4pt}
  \end{center}
  \begin{center}
  {\small Found a mistake or want to contribute a sheet?}\\
  {\small Visit the open-source project at \href{https://github.com/The-CerealDev/speed-maths}{github.com/The-CerealDev/speed-maths}}
  \end{center}
}

% Editing THIS file changes the look of every sheet/answer across every pillar.
% ══════════════════════════════════════════════════════════════════════════════

\usepackage[a4paper, top=25mm, bottom=25mm, left=22mm, right=22mm]{geometry}
\usepackage{amsmath, amssymb}
\usepackage{enumitem}
\usepackage{booktabs}
\usepackage{array}
\usepackage{parskip}
\usepackage{microtype}
\usepackage{fancyhdr}
\usepackage{titlesec}
\usepackage{mdframed}
\usepackage{xcolor}
\usepackage[hidelinks]{hyperref}
\usepackage{graphicx}
\usepackage{tikz}
\usepackage{pgfplots}
\pgfplotsset{compat=1.18}

% ── Colours ───────────────────────────────────────────────────────────────────
\definecolor{darkgray}{gray}{0.15}
\definecolor{medgray}{gray}{0.45}
\definecolor{lightgray}{gray}{0.93}
\definecolor{boxgray}{gray}{0.88}

% ── Section formatting ──────────────────────────────────────────────────────────
\titleformat{\section}[block]
  {\normalfont\large\bfseries}{}
  {0pt}{}[\vspace{2pt}\hrule\vspace{6pt}]
\titlespacing*{\section}{0pt}{14pt}{4pt}

% ── List formatting ─────────────────────────────────────────────────────────────
\setlist[enumerate,1]{label=\textbf{\arabic*.}, leftmargin=2em, itemsep=10pt}

% ── Branding constants (edit here, not per-file) ────────────────────────────────
\newcommand{\SpeedExamLine}{\textbf{Speed Maths:} Competition \& Exam Prep}
\newcommand{\SpeedCredit}{Series by \href{https://www.linkedin.com/in/cerealdev/}{David Akinyele-Aje}}

% ── Header / footer ──────────────────────────────────────────────────────────────
% Usage (once, after % main (standalone preamble for editing)
% vim: nowrap: spell:
% ══════════════════════════════════════════════════════════════════════════════
% Speed Maths --- shared preamble
% Included by every sheet and answer file via:  \input{../../shared/preamble}
% Editing THIS file changes the look of every sheet/answer across every pillar.
% ══════════════════════════════════════════════════════════════════════════════

\usepackage[a4paper, top=25mm, bottom=25mm, left=22mm, right=22mm]{geometry}
\usepackage{amsmath, amssymb}
\usepackage{enumitem}
\usepackage{booktabs}
\usepackage{array}
\usepackage{parskip}
\usepackage{microtype}
\usepackage{fancyhdr}
\usepackage{titlesec}
\usepackage{mdframed}
\usepackage{xcolor}
\usepackage[hidelinks]{hyperref}
\usepackage{graphicx}
\usepackage{tikz}
\usepackage{pgfplots}
\pgfplotsset{compat=1.18}

% ── Colours ───────────────────────────────────────────────────────────────────
\definecolor{darkgray}{gray}{0.15}
\definecolor{medgray}{gray}{0.45}
\definecolor{lightgray}{gray}{0.93}
\definecolor{boxgray}{gray}{0.88}

% ── Section formatting ──────────────────────────────────────────────────────────
\titleformat{\section}[block]
  {\normalfont\large\bfseries}{}
  {0pt}{}[\vspace{2pt}\hrule\vspace{6pt}]
\titlespacing*{\section}{0pt}{14pt}{4pt}

% ── List formatting ─────────────────────────────────────────────────────────────
\setlist[enumerate,1]{label=\textbf{\arabic*.}, leftmargin=2em, itemsep=10pt}

% ── Branding constants (edit here, not per-file) ────────────────────────────────
\newcommand{\SpeedExamLine}{\textbf{Speed Maths:} Competition \& Exam Prep}
\newcommand{\SpeedCredit}{Series by \href{https://www.linkedin.com/in/cerealdev/}{David Akinyele-Aje}}

% ── Header / footer ──────────────────────────────────────────────────────────────
% Usage (once, after \input{../../shared/preamble} and before \begin{document}):
%   \SpeedHeader{Algebra}{3}
% First arg  = pillar name shown as "Speed <Pillar> --- Daily Drill"
% Second arg = worksheet number
\pagestyle{fancy}
\newcommand{\SpeedHeader}[2]{%
  \fancyhf{}%
  \renewcommand{\headrulewidth}{0.4pt}%
  \setlength{\headheight}{14pt}%
  \fancyhead[L]{\small\textbf{Speed #1 --- Daily Drill}}%
  \fancyhead[R]{\small Worksheet \##2}%
  \fancyfoot[L]{\small \SpeedExamLine}%
  \fancyfoot[C]{\small\thepage}%
  \fancyfoot[R]{\small \SpeedCredit}%
  \renewcommand{\footrulewidth}{0.4pt}%
}

% ── Title block (used at the top of every sheet AND answer file) ───────────────
% Usage: \SpeedTitleBlock{Daily Algebra Drill \#3}{Jane Doe}
% %% AGENTS: When generating a new sheet, ask the human contributor for the name they want credited in argument 2.
% For answer files, pass the "--- Answers \& Investigations" suffix in the arg.
\newcommand{\SpeedTitleBlock}[2]{%
  \begin{center}
    {\LARGE\bfseries #1}\\[4pt]
    {\normalsize\itshape Sheet by #2}\\[6pt]
    \hrule\vspace{4pt}
    \begin{tabular}{llcc}
      \toprule
      \textbf{Section} & \textbf{Topic} & \textbf{Questions} & \textbf{Target time}\\
      \midrule
      A & Rapid Recognition          & 10 & 2 min 30 sec \\
      B & Manipulation Drills        & 10 & 8 min        \\
      C & Substitution \& Structure  & 8  & 10 min       \\
      D & Challenge                  & 5  & 15 min       \\
      \bottomrule
    \end{tabular}
    \vspace{4pt}
    \hrule
  \end{center}
}

% ── Sheet metadata: topic + the day's new tools ────────────────────────────────
% Usage (immediately after \SpeedTitleBlock, sheets only):
%   \SpeedMeta{Expanding, factorising, and the identity toolkit}
%             {difference of squares, sum/product form, completing the square}
%
% Arg 1 = the sheet's topic in one line. The website lists this instead of
%         "Sheet 03", so write the thing a reader would actually search for.
% Arg 2 = comma-separated, at most five: the tools this sheet introduces that
%         earlier sheets in the pillar did not. These become the website's tags.
%         Leave out anything the reader already met — the tags are meant to say
%         what is new today, not to summarise the sheet.
%
% Both are parsed straight out of this file by tools/build_website.py, so the
% sheet stays the single source of truth and the site cannot drift from it.
%
% This renders NOTHING. It is a declaration for the build script, not a piece of
% the sheet's layout: adding it to a sheet must not change that sheet's PDF, so
% the 25 existing sheets can be annotated without a single one of them being
% reissued. If the toolkit should also appear on the page, that is a separate
% decision about the sheet design — make it deliberately, here, not by leaving
% this macro printing something nobody asked it to.
\newcommand{\SpeedMeta}[2]{}

% ── Answer / Method / Investigate-further boxes (answer files only) ────────────
\newmdenv[
  backgroundcolor=lightgray,
  linecolor=medgray,
  linewidth=0.5pt,
  leftmargin=0pt, rightmargin=0pt,
  innerleftmargin=8pt, innerrightmargin=8pt,
  innertopmargin=5pt, innerbottommargin=5pt,
  skipabove=4pt, skipbelow=4pt
]{ansbox}

\newmdenv[
  backgroundcolor=white,
  linecolor=darkgray,
  linewidth=0.8pt,
  leftline=true, rightline=false, topline=false, bottomline=false,
  leftmargin=0pt, rightmargin=0pt,
  innerleftmargin=8pt, innerrightmargin=6pt,
  innertopmargin=4pt, innerbottommargin=4pt,
  skipabove=4pt, skipbelow=6pt
]{invbox}

\newcommand{\ans}[1]{%
  \begin{ansbox}
    \textbf{Answer:} #1
  \end{ansbox}}

\newcommand{\method}[1]{%
  {\small\textbf{Method:} #1}\par}

\newcommand{\inv}[1]{%
  \begin{invbox}
    \textbf{\small Investigate further:} {\small #1}
  \end{invbox}}

% ── Misc helpers ────────────────────────────────────────────────────────────────
\newcommand{\qn}[1]{\textbf{#1.}\;}

% Mistake-linking: attach a Top 5 Patterns / Common Traps bullet to the question
% label(s) in this sheet where it actually shows up, e.g. \seealso{B6, D2}
\newcommand{\seealso}[1]{\hfill\textit{\footnotesize(see #1)}}

% ── Graph options macro ─────────────────────────────────────────────────────────
% Usage: \GraphOptions{tikz code for A}{tikz code for B}{tikz code for C}{tikz code for D}
\newcommand{\GraphOptions}[4]{%
  \begin{center}
    \begin{tabular}{cc}
      \textbf{A)} & \textbf{B)} \\
      #1 & #2 \\[10pt]
      \textbf{C)} & \textbf{D)} \\
      #3 & #4
    \end{tabular}
  \end{center}
}

% ── Closing motivational rule + cross-promo footer (sheets only) ────────────────
% Usage: \SpeedClosing{"Quote text."}
\newcommand{\SpeedClosing}[1]{%
  \vspace{20pt}
  \begin{center}
  \rule{0.6\textwidth}{0.4pt}\\[4pt]
  {\small\itshape #1}\\[4pt]
  \rule{0.6\textwidth}{0.4pt}
  \end{center}
  \begin{center}
  {\small Found a mistake or want to contribute a sheet?}\\
  {\small Visit the open-source project at \href{https://github.com/The-CerealDev/speed-maths}{github.com/The-CerealDev/speed-maths}}
  \end{center}
}
 and before \begin{document}):
%   \SpeedHeader{Algebra}{3}
% First arg  = pillar name shown as "Speed <Pillar> --- Daily Drill"
% Second arg = worksheet number
\pagestyle{fancy}
\newcommand{\SpeedHeader}[2]{%
  \fancyhf{}%
  \renewcommand{\headrulewidth}{0.4pt}%
  \setlength{\headheight}{14pt}%
  \fancyhead[L]{\small\textbf{Speed #1 --- Daily Drill}}%
  \fancyhead[R]{\small Worksheet \##2}%
  \fancyfoot[L]{\small \SpeedExamLine}%
  \fancyfoot[C]{\small\thepage}%
  \fancyfoot[R]{\small \SpeedCredit}%
  \renewcommand{\footrulewidth}{0.4pt}%
}

% ── Title block (used at the top of every sheet AND answer file) ───────────────
% Usage: \SpeedTitleBlock{Daily Algebra Drill \#3}{Jane Doe}
% %% AGENTS: When generating a new sheet, ask the human contributor for the name they want credited in argument 2.
% For answer files, pass the "--- Answers \& Investigations" suffix in the arg.
\newcommand{\SpeedTitleBlock}[2]{%
  \begin{center}
    {\LARGE\bfseries #1}\\[4pt]
    {\normalsize\itshape Sheet by #2}\\[6pt]
    \hrule\vspace{4pt}
    \begin{tabular}{llcc}
      \toprule
      \textbf{Section} & \textbf{Topic} & \textbf{Questions} & \textbf{Target time}\\
      \midrule
      A & Rapid Recognition          & 10 & 2 min 30 sec \\
      B & Manipulation Drills        & 10 & 8 min        \\
      C & Substitution \& Structure  & 8  & 10 min       \\
      D & Challenge                  & 5  & 15 min       \\
      \bottomrule
    \end{tabular}
    \vspace{4pt}
    \hrule
  \end{center}
}

% ── Sheet metadata: topic + the day's new tools ────────────────────────────────
% Usage (immediately after \SpeedTitleBlock, sheets only):
%   \SpeedMeta{Expanding, factorising, and the identity toolkit}
%             {difference of squares, sum/product form, completing the square}
%
% Arg 1 = the sheet's topic in one line. The website lists this instead of
%         "Sheet 03", so write the thing a reader would actually search for.
% Arg 2 = comma-separated, at most five: the tools this sheet introduces that
%         earlier sheets in the pillar did not. These become the website's tags.
%         Leave out anything the reader already met — the tags are meant to say
%         what is new today, not to summarise the sheet.
%
% Both are parsed straight out of this file by tools/build_website.py, so the
% sheet stays the single source of truth and the site cannot drift from it.
%
% This renders NOTHING. It is a declaration for the build script, not a piece of
% the sheet's layout: adding it to a sheet must not change that sheet's PDF, so
% the 25 existing sheets can be annotated without a single one of them being
% reissued. If the toolkit should also appear on the page, that is a separate
% decision about the sheet design — make it deliberately, here, not by leaving
% this macro printing something nobody asked it to.
\newcommand{\SpeedMeta}[2]{}

% ── Answer / Method / Investigate-further boxes (answer files only) ────────────
\newmdenv[
  backgroundcolor=lightgray,
  linecolor=medgray,
  linewidth=0.5pt,
  leftmargin=0pt, rightmargin=0pt,
  innerleftmargin=8pt, innerrightmargin=8pt,
  innertopmargin=5pt, innerbottommargin=5pt,
  skipabove=4pt, skipbelow=4pt
]{ansbox}

\newmdenv[
  backgroundcolor=white,
  linecolor=darkgray,
  linewidth=0.8pt,
  leftline=true, rightline=false, topline=false, bottomline=false,
  leftmargin=0pt, rightmargin=0pt,
  innerleftmargin=8pt, innerrightmargin=6pt,
  innertopmargin=4pt, innerbottommargin=4pt,
  skipabove=4pt, skipbelow=6pt
]{invbox}

\newcommand{\ans}[1]{%
  \begin{ansbox}
    \textbf{Answer:} #1
  \end{ansbox}}

\newcommand{\method}[1]{%
  {\small\textbf{Method:} #1}\par}

\newcommand{\inv}[1]{%
  \begin{invbox}
    \textbf{\small Investigate further:} {\small #1}
  \end{invbox}}

% ── Misc helpers ────────────────────────────────────────────────────────────────
\newcommand{\qn}[1]{\textbf{#1.}\;}

% Mistake-linking: attach a Top 5 Patterns / Common Traps bullet to the question
% label(s) in this sheet where it actually shows up, e.g. \seealso{B6, D2}
\newcommand{\seealso}[1]{\hfill\textit{\footnotesize(see #1)}}

% ── Graph options macro ─────────────────────────────────────────────────────────
% Usage: \GraphOptions{tikz code for A}{tikz code for B}{tikz code for C}{tikz code for D}
\newcommand{\GraphOptions}[4]{%
  \begin{center}
    \begin{tabular}{cc}
      \textbf{A)} & \textbf{B)} \\
      #1 & #2 \\[10pt]
      \textbf{C)} & \textbf{D)} \\
      #3 & #4
    \end{tabular}
  \end{center}
}

% ── Closing motivational rule + cross-promo footer (sheets only) ────────────────
% Usage: \SpeedClosing{"Quote text."}
\newcommand{\SpeedClosing}[1]{%
  \vspace{20pt}
  \begin{center}
  \rule{0.6\textwidth}{0.4pt}\\[4pt]
  {\small\itshape #1}\\[4pt]
  \rule{0.6\textwidth}{0.4pt}
  \end{center}
  \begin{center}
  {\small Found a mistake or want to contribute a sheet?}\\
  {\small Visit the open-source project at \href{https://github.com/The-CerealDev/speed-maths}{github.com/The-CerealDev/speed-maths}}
  \end{center}
}

% Editing THIS file changes the look of every sheet/answer across every pillar.
% ══════════════════════════════════════════════════════════════════════════════

\usepackage[a4paper, top=25mm, bottom=25mm, left=22mm, right=22mm]{geometry}
\usepackage{amsmath, amssymb}
\usepackage{enumitem}
\usepackage{booktabs}
\usepackage{array}
\usepackage{parskip}
\usepackage{microtype}
\usepackage{fancyhdr}
\usepackage{titlesec}
\usepackage{mdframed}
\usepackage{xcolor}
\usepackage[hidelinks]{hyperref}
\usepackage{graphicx}
\usepackage{tikz}
\usepackage{pgfplots}
\pgfplotsset{compat=1.18}

% ── Colours ───────────────────────────────────────────────────────────────────
\definecolor{darkgray}{gray}{0.15}
\definecolor{medgray}{gray}{0.45}
\definecolor{lightgray}{gray}{0.93}
\definecolor{boxgray}{gray}{0.88}

% ── Section formatting ──────────────────────────────────────────────────────────
\titleformat{\section}[block]
  {\normalfont\large\bfseries}{}
  {0pt}{}[\vspace{2pt}\hrule\vspace{6pt}]
\titlespacing*{\section}{0pt}{14pt}{4pt}

% ── List formatting ─────────────────────────────────────────────────────────────
\setlist[enumerate,1]{label=\textbf{\arabic*.}, leftmargin=2em, itemsep=10pt}

% ── Branding constants (edit here, not per-file) ────────────────────────────────
\newcommand{\SpeedExamLine}{\textbf{Speed Maths:} Competition \& Exam Prep}
\newcommand{\SpeedCredit}{Series by \href{https://www.linkedin.com/in/cerealdev/}{David Akinyele-Aje}}

% ── Header / footer ──────────────────────────────────────────────────────────────
% Usage (once, after % main (standalone preamble for editing)
% vim: nowrap: spell:
% ══════════════════════════════════════════════════════════════════════════════
% Speed Maths --- shared preamble
% Included by every sheet and answer file via:  % main (standalone preamble for editing)
% vim: nowrap: spell:
% ══════════════════════════════════════════════════════════════════════════════
% Speed Maths --- shared preamble
% Included by every sheet and answer file via:  \input{../../shared/preamble}
% Editing THIS file changes the look of every sheet/answer across every pillar.
% ══════════════════════════════════════════════════════════════════════════════

\usepackage[a4paper, top=25mm, bottom=25mm, left=22mm, right=22mm]{geometry}
\usepackage{amsmath, amssymb}
\usepackage{enumitem}
\usepackage{booktabs}
\usepackage{array}
\usepackage{parskip}
\usepackage{microtype}
\usepackage{fancyhdr}
\usepackage{titlesec}
\usepackage{mdframed}
\usepackage{xcolor}
\usepackage[hidelinks]{hyperref}
\usepackage{graphicx}
\usepackage{tikz}
\usepackage{pgfplots}
\pgfplotsset{compat=1.18}

% ── Colours ───────────────────────────────────────────────────────────────────
\definecolor{darkgray}{gray}{0.15}
\definecolor{medgray}{gray}{0.45}
\definecolor{lightgray}{gray}{0.93}
\definecolor{boxgray}{gray}{0.88}

% ── Section formatting ──────────────────────────────────────────────────────────
\titleformat{\section}[block]
  {\normalfont\large\bfseries}{}
  {0pt}{}[\vspace{2pt}\hrule\vspace{6pt}]
\titlespacing*{\section}{0pt}{14pt}{4pt}

% ── List formatting ─────────────────────────────────────────────────────────────
\setlist[enumerate,1]{label=\textbf{\arabic*.}, leftmargin=2em, itemsep=10pt}

% ── Branding constants (edit here, not per-file) ────────────────────────────────
\newcommand{\SpeedExamLine}{\textbf{Speed Maths:} Competition \& Exam Prep}
\newcommand{\SpeedCredit}{Series by \href{https://www.linkedin.com/in/cerealdev/}{David Akinyele-Aje}}

% ── Header / footer ──────────────────────────────────────────────────────────────
% Usage (once, after \input{../../shared/preamble} and before \begin{document}):
%   \SpeedHeader{Algebra}{3}
% First arg  = pillar name shown as "Speed <Pillar> --- Daily Drill"
% Second arg = worksheet number
\pagestyle{fancy}
\newcommand{\SpeedHeader}[2]{%
  \fancyhf{}%
  \renewcommand{\headrulewidth}{0.4pt}%
  \setlength{\headheight}{14pt}%
  \fancyhead[L]{\small\textbf{Speed #1 --- Daily Drill}}%
  \fancyhead[R]{\small Worksheet \##2}%
  \fancyfoot[L]{\small \SpeedExamLine}%
  \fancyfoot[C]{\small\thepage}%
  \fancyfoot[R]{\small \SpeedCredit}%
  \renewcommand{\footrulewidth}{0.4pt}%
}

% ── Title block (used at the top of every sheet AND answer file) ───────────────
% Usage: \SpeedTitleBlock{Daily Algebra Drill \#3}{Jane Doe}
% %% AGENTS: When generating a new sheet, ask the human contributor for the name they want credited in argument 2.
% For answer files, pass the "--- Answers \& Investigations" suffix in the arg.
\newcommand{\SpeedTitleBlock}[2]{%
  \begin{center}
    {\LARGE\bfseries #1}\\[4pt]
    {\normalsize\itshape Sheet by #2}\\[6pt]
    \hrule\vspace{4pt}
    \begin{tabular}{llcc}
      \toprule
      \textbf{Section} & \textbf{Topic} & \textbf{Questions} & \textbf{Target time}\\
      \midrule
      A & Rapid Recognition          & 10 & 2 min 30 sec \\
      B & Manipulation Drills        & 10 & 8 min        \\
      C & Substitution \& Structure  & 8  & 10 min       \\
      D & Challenge                  & 5  & 15 min       \\
      \bottomrule
    \end{tabular}
    \vspace{4pt}
    \hrule
  \end{center}
}

% ── Sheet metadata: topic + the day's new tools ────────────────────────────────
% Usage (immediately after \SpeedTitleBlock, sheets only):
%   \SpeedMeta{Expanding, factorising, and the identity toolkit}
%             {difference of squares, sum/product form, completing the square}
%
% Arg 1 = the sheet's topic in one line. The website lists this instead of
%         "Sheet 03", so write the thing a reader would actually search for.
% Arg 2 = comma-separated, at most five: the tools this sheet introduces that
%         earlier sheets in the pillar did not. These become the website's tags.
%         Leave out anything the reader already met — the tags are meant to say
%         what is new today, not to summarise the sheet.
%
% Both are parsed straight out of this file by tools/build_website.py, so the
% sheet stays the single source of truth and the site cannot drift from it.
%
% This renders NOTHING. It is a declaration for the build script, not a piece of
% the sheet's layout: adding it to a sheet must not change that sheet's PDF, so
% the 25 existing sheets can be annotated without a single one of them being
% reissued. If the toolkit should also appear on the page, that is a separate
% decision about the sheet design — make it deliberately, here, not by leaving
% this macro printing something nobody asked it to.
\newcommand{\SpeedMeta}[2]{}

% ── Answer / Method / Investigate-further boxes (answer files only) ────────────
\newmdenv[
  backgroundcolor=lightgray,
  linecolor=medgray,
  linewidth=0.5pt,
  leftmargin=0pt, rightmargin=0pt,
  innerleftmargin=8pt, innerrightmargin=8pt,
  innertopmargin=5pt, innerbottommargin=5pt,
  skipabove=4pt, skipbelow=4pt
]{ansbox}

\newmdenv[
  backgroundcolor=white,
  linecolor=darkgray,
  linewidth=0.8pt,
  leftline=true, rightline=false, topline=false, bottomline=false,
  leftmargin=0pt, rightmargin=0pt,
  innerleftmargin=8pt, innerrightmargin=6pt,
  innertopmargin=4pt, innerbottommargin=4pt,
  skipabove=4pt, skipbelow=6pt
]{invbox}

\newcommand{\ans}[1]{%
  \begin{ansbox}
    \textbf{Answer:} #1
  \end{ansbox}}

\newcommand{\method}[1]{%
  {\small\textbf{Method:} #1}\par}

\newcommand{\inv}[1]{%
  \begin{invbox}
    \textbf{\small Investigate further:} {\small #1}
  \end{invbox}}

% ── Misc helpers ────────────────────────────────────────────────────────────────
\newcommand{\qn}[1]{\textbf{#1.}\;}

% Mistake-linking: attach a Top 5 Patterns / Common Traps bullet to the question
% label(s) in this sheet where it actually shows up, e.g. \seealso{B6, D2}
\newcommand{\seealso}[1]{\hfill\textit{\footnotesize(see #1)}}

% ── Graph options macro ─────────────────────────────────────────────────────────
% Usage: \GraphOptions{tikz code for A}{tikz code for B}{tikz code for C}{tikz code for D}
\newcommand{\GraphOptions}[4]{%
  \begin{center}
    \begin{tabular}{cc}
      \textbf{A)} & \textbf{B)} \\
      #1 & #2 \\[10pt]
      \textbf{C)} & \textbf{D)} \\
      #3 & #4
    \end{tabular}
  \end{center}
}

% ── Closing motivational rule + cross-promo footer (sheets only) ────────────────
% Usage: \SpeedClosing{"Quote text."}
\newcommand{\SpeedClosing}[1]{%
  \vspace{20pt}
  \begin{center}
  \rule{0.6\textwidth}{0.4pt}\\[4pt]
  {\small\itshape #1}\\[4pt]
  \rule{0.6\textwidth}{0.4pt}
  \end{center}
  \begin{center}
  {\small Found a mistake or want to contribute a sheet?}\\
  {\small Visit the open-source project at \href{https://github.com/The-CerealDev/speed-maths}{github.com/The-CerealDev/speed-maths}}
  \end{center}
}

% Editing THIS file changes the look of every sheet/answer across every pillar.
% ══════════════════════════════════════════════════════════════════════════════

\usepackage[a4paper, top=25mm, bottom=25mm, left=22mm, right=22mm]{geometry}
\usepackage{amsmath, amssymb}
\usepackage{enumitem}
\usepackage{booktabs}
\usepackage{array}
\usepackage{parskip}
\usepackage{microtype}
\usepackage{fancyhdr}
\usepackage{titlesec}
\usepackage{mdframed}
\usepackage{xcolor}
\usepackage[hidelinks]{hyperref}
\usepackage{graphicx}
\usepackage{tikz}
\usepackage{pgfplots}
\pgfplotsset{compat=1.18}

% ── Colours ───────────────────────────────────────────────────────────────────
\definecolor{darkgray}{gray}{0.15}
\definecolor{medgray}{gray}{0.45}
\definecolor{lightgray}{gray}{0.93}
\definecolor{boxgray}{gray}{0.88}

% ── Section formatting ──────────────────────────────────────────────────────────
\titleformat{\section}[block]
  {\normalfont\large\bfseries}{}
  {0pt}{}[\vspace{2pt}\hrule\vspace{6pt}]
\titlespacing*{\section}{0pt}{14pt}{4pt}

% ── List formatting ─────────────────────────────────────────────────────────────
\setlist[enumerate,1]{label=\textbf{\arabic*.}, leftmargin=2em, itemsep=10pt}

% ── Branding constants (edit here, not per-file) ────────────────────────────────
\newcommand{\SpeedExamLine}{\textbf{Speed Maths:} Competition \& Exam Prep}
\newcommand{\SpeedCredit}{Series by \href{https://www.linkedin.com/in/cerealdev/}{David Akinyele-Aje}}

% ── Header / footer ──────────────────────────────────────────────────────────────
% Usage (once, after % main (standalone preamble for editing)
% vim: nowrap: spell:
% ══════════════════════════════════════════════════════════════════════════════
% Speed Maths --- shared preamble
% Included by every sheet and answer file via:  \input{../../shared/preamble}
% Editing THIS file changes the look of every sheet/answer across every pillar.
% ══════════════════════════════════════════════════════════════════════════════

\usepackage[a4paper, top=25mm, bottom=25mm, left=22mm, right=22mm]{geometry}
\usepackage{amsmath, amssymb}
\usepackage{enumitem}
\usepackage{booktabs}
\usepackage{array}
\usepackage{parskip}
\usepackage{microtype}
\usepackage{fancyhdr}
\usepackage{titlesec}
\usepackage{mdframed}
\usepackage{xcolor}
\usepackage[hidelinks]{hyperref}
\usepackage{graphicx}
\usepackage{tikz}
\usepackage{pgfplots}
\pgfplotsset{compat=1.18}

% ── Colours ───────────────────────────────────────────────────────────────────
\definecolor{darkgray}{gray}{0.15}
\definecolor{medgray}{gray}{0.45}
\definecolor{lightgray}{gray}{0.93}
\definecolor{boxgray}{gray}{0.88}

% ── Section formatting ──────────────────────────────────────────────────────────
\titleformat{\section}[block]
  {\normalfont\large\bfseries}{}
  {0pt}{}[\vspace{2pt}\hrule\vspace{6pt}]
\titlespacing*{\section}{0pt}{14pt}{4pt}

% ── List formatting ─────────────────────────────────────────────────────────────
\setlist[enumerate,1]{label=\textbf{\arabic*.}, leftmargin=2em, itemsep=10pt}

% ── Branding constants (edit here, not per-file) ────────────────────────────────
\newcommand{\SpeedExamLine}{\textbf{Speed Maths:} Competition \& Exam Prep}
\newcommand{\SpeedCredit}{Series by \href{https://www.linkedin.com/in/cerealdev/}{David Akinyele-Aje}}

% ── Header / footer ──────────────────────────────────────────────────────────────
% Usage (once, after \input{../../shared/preamble} and before \begin{document}):
%   \SpeedHeader{Algebra}{3}
% First arg  = pillar name shown as "Speed <Pillar> --- Daily Drill"
% Second arg = worksheet number
\pagestyle{fancy}
\newcommand{\SpeedHeader}[2]{%
  \fancyhf{}%
  \renewcommand{\headrulewidth}{0.4pt}%
  \setlength{\headheight}{14pt}%
  \fancyhead[L]{\small\textbf{Speed #1 --- Daily Drill}}%
  \fancyhead[R]{\small Worksheet \##2}%
  \fancyfoot[L]{\small \SpeedExamLine}%
  \fancyfoot[C]{\small\thepage}%
  \fancyfoot[R]{\small \SpeedCredit}%
  \renewcommand{\footrulewidth}{0.4pt}%
}

% ── Title block (used at the top of every sheet AND answer file) ───────────────
% Usage: \SpeedTitleBlock{Daily Algebra Drill \#3}{Jane Doe}
% %% AGENTS: When generating a new sheet, ask the human contributor for the name they want credited in argument 2.
% For answer files, pass the "--- Answers \& Investigations" suffix in the arg.
\newcommand{\SpeedTitleBlock}[2]{%
  \begin{center}
    {\LARGE\bfseries #1}\\[4pt]
    {\normalsize\itshape Sheet by #2}\\[6pt]
    \hrule\vspace{4pt}
    \begin{tabular}{llcc}
      \toprule
      \textbf{Section} & \textbf{Topic} & \textbf{Questions} & \textbf{Target time}\\
      \midrule
      A & Rapid Recognition          & 10 & 2 min 30 sec \\
      B & Manipulation Drills        & 10 & 8 min        \\
      C & Substitution \& Structure  & 8  & 10 min       \\
      D & Challenge                  & 5  & 15 min       \\
      \bottomrule
    \end{tabular}
    \vspace{4pt}
    \hrule
  \end{center}
}

% ── Sheet metadata: topic + the day's new tools ────────────────────────────────
% Usage (immediately after \SpeedTitleBlock, sheets only):
%   \SpeedMeta{Expanding, factorising, and the identity toolkit}
%             {difference of squares, sum/product form, completing the square}
%
% Arg 1 = the sheet's topic in one line. The website lists this instead of
%         "Sheet 03", so write the thing a reader would actually search for.
% Arg 2 = comma-separated, at most five: the tools this sheet introduces that
%         earlier sheets in the pillar did not. These become the website's tags.
%         Leave out anything the reader already met — the tags are meant to say
%         what is new today, not to summarise the sheet.
%
% Both are parsed straight out of this file by tools/build_website.py, so the
% sheet stays the single source of truth and the site cannot drift from it.
%
% This renders NOTHING. It is a declaration for the build script, not a piece of
% the sheet's layout: adding it to a sheet must not change that sheet's PDF, so
% the 25 existing sheets can be annotated without a single one of them being
% reissued. If the toolkit should also appear on the page, that is a separate
% decision about the sheet design — make it deliberately, here, not by leaving
% this macro printing something nobody asked it to.
\newcommand{\SpeedMeta}[2]{}

% ── Answer / Method / Investigate-further boxes (answer files only) ────────────
\newmdenv[
  backgroundcolor=lightgray,
  linecolor=medgray,
  linewidth=0.5pt,
  leftmargin=0pt, rightmargin=0pt,
  innerleftmargin=8pt, innerrightmargin=8pt,
  innertopmargin=5pt, innerbottommargin=5pt,
  skipabove=4pt, skipbelow=4pt
]{ansbox}

\newmdenv[
  backgroundcolor=white,
  linecolor=darkgray,
  linewidth=0.8pt,
  leftline=true, rightline=false, topline=false, bottomline=false,
  leftmargin=0pt, rightmargin=0pt,
  innerleftmargin=8pt, innerrightmargin=6pt,
  innertopmargin=4pt, innerbottommargin=4pt,
  skipabove=4pt, skipbelow=6pt
]{invbox}

\newcommand{\ans}[1]{%
  \begin{ansbox}
    \textbf{Answer:} #1
  \end{ansbox}}

\newcommand{\method}[1]{%
  {\small\textbf{Method:} #1}\par}

\newcommand{\inv}[1]{%
  \begin{invbox}
    \textbf{\small Investigate further:} {\small #1}
  \end{invbox}}

% ── Misc helpers ────────────────────────────────────────────────────────────────
\newcommand{\qn}[1]{\textbf{#1.}\;}

% Mistake-linking: attach a Top 5 Patterns / Common Traps bullet to the question
% label(s) in this sheet where it actually shows up, e.g. \seealso{B6, D2}
\newcommand{\seealso}[1]{\hfill\textit{\footnotesize(see #1)}}

% ── Graph options macro ─────────────────────────────────────────────────────────
% Usage: \GraphOptions{tikz code for A}{tikz code for B}{tikz code for C}{tikz code for D}
\newcommand{\GraphOptions}[4]{%
  \begin{center}
    \begin{tabular}{cc}
      \textbf{A)} & \textbf{B)} \\
      #1 & #2 \\[10pt]
      \textbf{C)} & \textbf{D)} \\
      #3 & #4
    \end{tabular}
  \end{center}
}

% ── Closing motivational rule + cross-promo footer (sheets only) ────────────────
% Usage: \SpeedClosing{"Quote text."}
\newcommand{\SpeedClosing}[1]{%
  \vspace{20pt}
  \begin{center}
  \rule{0.6\textwidth}{0.4pt}\\[4pt]
  {\small\itshape #1}\\[4pt]
  \rule{0.6\textwidth}{0.4pt}
  \end{center}
  \begin{center}
  {\small Found a mistake or want to contribute a sheet?}\\
  {\small Visit the open-source project at \href{https://github.com/The-CerealDev/speed-maths}{github.com/The-CerealDev/speed-maths}}
  \end{center}
}
 and before \begin{document}):
%   \SpeedHeader{Algebra}{3}
% First arg  = pillar name shown as "Speed <Pillar> --- Daily Drill"
% Second arg = worksheet number
\pagestyle{fancy}
\newcommand{\SpeedHeader}[2]{%
  \fancyhf{}%
  \renewcommand{\headrulewidth}{0.4pt}%
  \setlength{\headheight}{14pt}%
  \fancyhead[L]{\small\textbf{Speed #1 --- Daily Drill}}%
  \fancyhead[R]{\small Worksheet \##2}%
  \fancyfoot[L]{\small \SpeedExamLine}%
  \fancyfoot[C]{\small\thepage}%
  \fancyfoot[R]{\small \SpeedCredit}%
  \renewcommand{\footrulewidth}{0.4pt}%
}

% ── Title block (used at the top of every sheet AND answer file) ───────────────
% Usage: \SpeedTitleBlock{Daily Algebra Drill \#3}{Jane Doe}
% %% AGENTS: When generating a new sheet, ask the human contributor for the name they want credited in argument 2.
% For answer files, pass the "--- Answers \& Investigations" suffix in the arg.
\newcommand{\SpeedTitleBlock}[2]{%
  \begin{center}
    {\LARGE\bfseries #1}\\[4pt]
    {\normalsize\itshape Sheet by #2}\\[6pt]
    \hrule\vspace{4pt}
    \begin{tabular}{llcc}
      \toprule
      \textbf{Section} & \textbf{Topic} & \textbf{Questions} & \textbf{Target time}\\
      \midrule
      A & Rapid Recognition          & 10 & 2 min 30 sec \\
      B & Manipulation Drills        & 10 & 8 min        \\
      C & Substitution \& Structure  & 8  & 10 min       \\
      D & Challenge                  & 5  & 15 min       \\
      \bottomrule
    \end{tabular}
    \vspace{4pt}
    \hrule
  \end{center}
}

% ── Sheet metadata: topic + the day's new tools ────────────────────────────────
% Usage (immediately after \SpeedTitleBlock, sheets only):
%   \SpeedMeta{Expanding, factorising, and the identity toolkit}
%             {difference of squares, sum/product form, completing the square}
%
% Arg 1 = the sheet's topic in one line. The website lists this instead of
%         "Sheet 03", so write the thing a reader would actually search for.
% Arg 2 = comma-separated, at most five: the tools this sheet introduces that
%         earlier sheets in the pillar did not. These become the website's tags.
%         Leave out anything the reader already met — the tags are meant to say
%         what is new today, not to summarise the sheet.
%
% Both are parsed straight out of this file by tools/build_website.py, so the
% sheet stays the single source of truth and the site cannot drift from it.
%
% This renders NOTHING. It is a declaration for the build script, not a piece of
% the sheet's layout: adding it to a sheet must not change that sheet's PDF, so
% the 25 existing sheets can be annotated without a single one of them being
% reissued. If the toolkit should also appear on the page, that is a separate
% decision about the sheet design — make it deliberately, here, not by leaving
% this macro printing something nobody asked it to.
\newcommand{\SpeedMeta}[2]{}

% ── Answer / Method / Investigate-further boxes (answer files only) ────────────
\newmdenv[
  backgroundcolor=lightgray,
  linecolor=medgray,
  linewidth=0.5pt,
  leftmargin=0pt, rightmargin=0pt,
  innerleftmargin=8pt, innerrightmargin=8pt,
  innertopmargin=5pt, innerbottommargin=5pt,
  skipabove=4pt, skipbelow=4pt
]{ansbox}

\newmdenv[
  backgroundcolor=white,
  linecolor=darkgray,
  linewidth=0.8pt,
  leftline=true, rightline=false, topline=false, bottomline=false,
  leftmargin=0pt, rightmargin=0pt,
  innerleftmargin=8pt, innerrightmargin=6pt,
  innertopmargin=4pt, innerbottommargin=4pt,
  skipabove=4pt, skipbelow=6pt
]{invbox}

\newcommand{\ans}[1]{%
  \begin{ansbox}
    \textbf{Answer:} #1
  \end{ansbox}}

\newcommand{\method}[1]{%
  {\small\textbf{Method:} #1}\par}

\newcommand{\inv}[1]{%
  \begin{invbox}
    \textbf{\small Investigate further:} {\small #1}
  \end{invbox}}

% ── Misc helpers ────────────────────────────────────────────────────────────────
\newcommand{\qn}[1]{\textbf{#1.}\;}

% Mistake-linking: attach a Top 5 Patterns / Common Traps bullet to the question
% label(s) in this sheet where it actually shows up, e.g. \seealso{B6, D2}
\newcommand{\seealso}[1]{\hfill\textit{\footnotesize(see #1)}}

% ── Graph options macro ─────────────────────────────────────────────────────────
% Usage: \GraphOptions{tikz code for A}{tikz code for B}{tikz code for C}{tikz code for D}
\newcommand{\GraphOptions}[4]{%
  \begin{center}
    \begin{tabular}{cc}
      \textbf{A)} & \textbf{B)} \\
      #1 & #2 \\[10pt]
      \textbf{C)} & \textbf{D)} \\
      #3 & #4
    \end{tabular}
  \end{center}
}

% ── Closing motivational rule + cross-promo footer (sheets only) ────────────────
% Usage: \SpeedClosing{"Quote text."}
\newcommand{\SpeedClosing}[1]{%
  \vspace{20pt}
  \begin{center}
  \rule{0.6\textwidth}{0.4pt}\\[4pt]
  {\small\itshape #1}\\[4pt]
  \rule{0.6\textwidth}{0.4pt}
  \end{center}
  \begin{center}
  {\small Found a mistake or want to contribute a sheet?}\\
  {\small Visit the open-source project at \href{https://github.com/The-CerealDev/speed-maths}{github.com/The-CerealDev/speed-maths}}
  \end{center}
}
 and before \begin{document}):
%   \SpeedHeader{Algebra}{3}
% First arg  = pillar name shown as "Speed <Pillar> --- Daily Drill"
% Second arg = worksheet number
\pagestyle{fancy}
\newcommand{\SpeedHeader}[2]{%
  \fancyhf{}%
  \renewcommand{\headrulewidth}{0.4pt}%
  \setlength{\headheight}{14pt}%
  \fancyhead[L]{\small\textbf{Speed #1 --- Daily Drill}}%
  \fancyhead[R]{\small Worksheet \##2}%
  \fancyfoot[L]{\small \SpeedExamLine}%
  \fancyfoot[C]{\small\thepage}%
  \fancyfoot[R]{\small \SpeedCredit}%
  \renewcommand{\footrulewidth}{0.4pt}%
}

% ── Title block (used at the top of every sheet AND answer file) ───────────────
% Usage: \SpeedTitleBlock{Daily Algebra Drill \#3}{Jane Doe}
% %% AGENTS: When generating a new sheet, ask the human contributor for the name they want credited in argument 2.
% For answer files, pass the "--- Answers \& Investigations" suffix in the arg.
\newcommand{\SpeedTitleBlock}[2]{%
  \begin{center}
    {\LARGE\bfseries #1}\\[4pt]
    {\normalsize\itshape Sheet by #2}\\[6pt]
    \hrule\vspace{4pt}
    \begin{tabular}{llcc}
      \toprule
      \textbf{Section} & \textbf{Topic} & \textbf{Questions} & \textbf{Target time}\\
      \midrule
      A & Rapid Recognition          & 10 & 2 min 30 sec \\
      B & Manipulation Drills        & 10 & 8 min        \\
      C & Substitution \& Structure  & 8  & 10 min       \\
      D & Challenge                  & 5  & 15 min       \\
      \bottomrule
    \end{tabular}
    \vspace{4pt}
    \hrule
  \end{center}
}

% ── Sheet metadata: topic + the day's new tools ────────────────────────────────
% Usage (immediately after \SpeedTitleBlock, sheets only):
%   \SpeedMeta{Expanding, factorising, and the identity toolkit}
%             {difference of squares, sum/product form, completing the square}
%
% Arg 1 = the sheet's topic in one line. The website lists this instead of
%         "Sheet 03", so write the thing a reader would actually search for.
% Arg 2 = comma-separated, at most five: the tools this sheet introduces that
%         earlier sheets in the pillar did not. These become the website's tags.
%         Leave out anything the reader already met — the tags are meant to say
%         what is new today, not to summarise the sheet.
%
% Both are parsed straight out of this file by tools/build_website.py, so the
% sheet stays the single source of truth and the site cannot drift from it.
%
% This renders NOTHING. It is a declaration for the build script, not a piece of
% the sheet's layout: adding it to a sheet must not change that sheet's PDF, so
% the 25 existing sheets can be annotated without a single one of them being
% reissued. If the toolkit should also appear on the page, that is a separate
% decision about the sheet design — make it deliberately, here, not by leaving
% this macro printing something nobody asked it to.
\newcommand{\SpeedMeta}[2]{}

% ── Answer / Method / Investigate-further boxes (answer files only) ────────────
\newmdenv[
  backgroundcolor=lightgray,
  linecolor=medgray,
  linewidth=0.5pt,
  leftmargin=0pt, rightmargin=0pt,
  innerleftmargin=8pt, innerrightmargin=8pt,
  innertopmargin=5pt, innerbottommargin=5pt,
  skipabove=4pt, skipbelow=4pt
]{ansbox}

\newmdenv[
  backgroundcolor=white,
  linecolor=darkgray,
  linewidth=0.8pt,
  leftline=true, rightline=false, topline=false, bottomline=false,
  leftmargin=0pt, rightmargin=0pt,
  innerleftmargin=8pt, innerrightmargin=6pt,
  innertopmargin=4pt, innerbottommargin=4pt,
  skipabove=4pt, skipbelow=6pt
]{invbox}

\newcommand{\ans}[1]{%
  \begin{ansbox}
    \textbf{Answer:} #1
  \end{ansbox}}

\newcommand{\method}[1]{%
  {\small\textbf{Method:} #1}\par}

\newcommand{\inv}[1]{%
  \begin{invbox}
    \textbf{\small Investigate further:} {\small #1}
  \end{invbox}}

% ── Misc helpers ────────────────────────────────────────────────────────────────
\newcommand{\qn}[1]{\textbf{#1.}\;}

% Mistake-linking: attach a Top 5 Patterns / Common Traps bullet to the question
% label(s) in this sheet where it actually shows up, e.g. \seealso{B6, D2}
\newcommand{\seealso}[1]{\hfill\textit{\footnotesize(see #1)}}

% ── Graph options macro ─────────────────────────────────────────────────────────
% Usage: \GraphOptions{tikz code for A}{tikz code for B}{tikz code for C}{tikz code for D}
\newcommand{\GraphOptions}[4]{%
  \begin{center}
    \begin{tabular}{cc}
      \textbf{A)} & \textbf{B)} \\
      #1 & #2 \\[10pt]
      \textbf{C)} & \textbf{D)} \\
      #3 & #4
    \end{tabular}
  \end{center}
}

% ── Closing motivational rule + cross-promo footer (sheets only) ────────────────
% Usage: \SpeedClosing{"Quote text."}
\newcommand{\SpeedClosing}[1]{%
  \vspace{20pt}
  \begin{center}
  \rule{0.6\textwidth}{0.4pt}\\[4pt]
  {\small\itshape #1}\\[4pt]
  \rule{0.6\textwidth}{0.4pt}
  \end{center}
  \begin{center}
  {\small Found a mistake or want to contribute a sheet?}\\
  {\small Visit the open-source project at \href{https://github.com/The-CerealDev/speed-maths}{github.com/The-CerealDev/speed-maths}}
  \end{center}
}

\SpeedHeader{Logic}{6}

\begin{document}

\SpeedTitleBlock{Answers: Logic Drill \#6}{David Akinyele-Aje}

\section*{Section A}

\begin{enumerate}

%── A1 ──────────────────────────────────────────────────────────────────────────
\item What are the two essential components required for a proof by mathematical induction?

\ans{A base case, and an inductive step.}
\method{The base case establishes $P(n_0)$ for the starting value; the inductive step proves $P(k)\implies P(k+1)$ for all relevant $k$. Together these imply $P(n)$ for every $n\geq n_0$.}
\inv{Why is neither component alone sufficient?}

%── A2 ──────────────────────────────────────────────────────────────────────────
\item To prove $P(n)$ for all integers $n\geq4$ by induction, what is the base case?

\ans{$P(4)$}
\method{The base case is always the smallest value in the claimed range.}
\inv{If the claim were instead ``for all $n\geq4$ except $n=5$'', would ordinary induction still apply directly?}

%── A3 ──────────────────────────────────────────────────────────────────────────
\item If $P(k)\implies P(k+1)$ is proved for all $k\geq1$, but no base case is ever established, has $P(n)$ been proved for any $n$?

\ans{No.}
\method{The inductive step alone is vacuous without a starting point --- it is entirely possible for $P(k)\implies P(k+1)$ to hold for every $k$ while $P(n)$ is false for every $n$ (e.g.\ if $P(n)$ is always false, $P(k)\implies P(k+1)$ holds vacuously).}
\inv{Give an example of a statement $P(n)$ that is false for every $n$ but for which $P(k)\implies P(k+1)$ holds for all $k$.}

%── A4 ──────────────────────────────────────────────────────────────────────────
\item In a proof of $\sum_{i=1}^n i=\frac{n(n+1)}{2}$ by induction, what statement exactly is assumed as the inductive hypothesis?

\ans{$\sum_{i=1}^k i=\frac{k(k+1)}{2}$, for some fixed (but arbitrary) $k$.}
\method{The inductive hypothesis assumes the formula holds at $n=k$, in order to derive it at $n=k+1$.}
\inv{Why must $k$ be treated as arbitrary rather than a specific number?}

%── A5 ──────────────────────────────────────────────────────────────────────────
\item True or false: ``Mathematical induction can be used to prove statements about every real number $x\geq1$.''

\ans{False}
\method{Induction proves statements indexed by the integers (or a well-ordered discrete set), stepping from $k$ to $k+1$. The real numbers between consecutive integers are never reached by this stepping process.}
\inv{What technique is typically used instead for statements about all real $x$ in an interval?}

%── A6 ──────────────────────────────────────────────────────────────────────────
\item To prove $2^n>n^2$ for all $n\geq5$ by induction, what is the smallest value of $n$ for which the base case must be verified?

\ans{$n=5$}
\method{Check: $2^5=32>25=5^2$, true. (Note $2^4=16=4^2$ and $2^3=8<9=3^2$, so the claim genuinely fails before $n=5$ --- this is why the range starts at $5$.)}
\inv{Verify that the claim fails at $n=4$ and $n=3$.}

%── A7 ──────────────────────────────────────────────────────────────────────────
\item What is the negation of the inductive-step statement ``For all $k\geq1$, $P(k)\implies P(k+1)$''?

\ans{There exists $k\geq1$ such that $P(k)$ is true and $P(k+1)$ is false.}
\method{$\neg(\forall k\, (P(k)\implies P(k+1))) \equiv \exists k\,\neg(P(k)\implies P(k+1)) \equiv \exists k\,(P(k)\land\neg P(k+1))$.}
\inv{Which earlier drill covered exactly this negation rule for implications?}

%── A8 ──────────────────────────────────────────────────────────────────────────
\item True or false: ``If $P(1)$ is true and $P(k)\implies P(k+2)$ holds for all $k\geq1$, then $P(n)$ is true for every positive integer $n$.''

\ans{False}
\method{$P(1)\implies P(3)\implies P(5)\implies\cdots$ only establishes $P(n)$ for odd $n$. $P(2)$ is never reached, since nothing in the chain starts from an even index.}
\inv{What additional single fact would need to be assumed to fix this and cover every positive integer?}

%── A9 ──────────────────────────────────────────────────────────────────────────
\item What is the smallest positive integer $n$ for which $n!>3^n$?

\ans{$n=7$}
\method{$6!=720$ vs $3^6=729$: $720<729$, fails. $7!=5040$ vs $3^7=2187$: $5040>2187$, holds. So $n=7$ is the smallest.}
\inv{Verify $n!\leq3^n$ for $n=1,\dots,6$ explicitly.}

%── A10 ──────────────────────────────────────────────────────────────────────────
\item In a proof by strong induction, what may the inductive step assume that ordinary (``weak'') induction's inductive step may not?

\ans{The truth of $P(1),P(2),\dots,P(k)$ (all previous cases), not just $P(k)$ alone.}
\method{Weak induction's inductive step assumes only $P(k)$ to derive $P(k+1)$; strong induction's inductive step may use any or all of $P(1),\dots,P(k)$.}
\inv{Are weak and strong induction actually logically equivalent in what they can prove, despite the difference in what's assumed?}

\end{enumerate}

\section*{Section B}

\begin{enumerate}

%── B1 ──────────────────────────────────────────────────────────────────────────
\item Express $\sum_{i=1}^{k+1} i^3$ in terms of $\sum_{i=1}^k i^3$.

\ans{$\sum_{i=1}^{k+1} i^3 = \sum_{i=1}^k i^3 + (k+1)^3$}
\method{The sum up to $k+1$ terms equals the sum up to $k$ terms plus the new $(k+1)$-th term.}
\inv{If $\sum_{i=1}^k i^3=\left(\frac{k(k+1)}2\right)^2$, verify the inductive step closes correctly.}

%── B2 ──────────────────────────────────────────────────────────────────────────
\item Assuming $\sum_{i=1}^k i=\frac{k(k+1)}{2}$, add $(k+1)$ and simplify to show $\sum_{i=1}^{k+1} i=\frac{(k+1)(k+2)}{2}$.

\ans{$\frac{k(k+1)}2+(k+1)=\frac{k(k+1)+2(k+1)}2=\frac{(k+1)(k+2)}2$}
\method{Factor out $(k+1)$ from both terms in the numerator.}
\inv{What is the base case for this induction, and does it check out?}

%── B3 ──────────────────────────────────────────────────────────────────────────
\item Prove by induction: $\sum_{i=1}^n (2i-1)=n^2$ for all positive integers $n$.

\ans{Proof by induction}
\method{Base case $n=1$: LHS $=1$, RHS $=1^2=1$. Holds.\\
Inductive step: assume $\sum_{i=1}^k(2i-1)=k^2$. Then $\sum_{i=1}^{k+1}(2i-1)=k^2+(2(k+1)-1)=k^2+2k+1=(k+1)^2$. So the formula holds for $k+1$. By induction, true for all positive integers $n$.}
\inv{What is the geometric interpretation of this identity (sum of first $n$ odd numbers)?}

%── B4 ──────────────────────────────────────────────────────────────────────────
\item Assuming $4^{2k}-1$ is divisible by $15$, express $4^{2(k+1)}-1$ in a form that shows it is also divisible by $15$.

\ans{$4^{2(k+1)}-1=16(4^{2k}-1)+15$}
\method{$4^{2(k+1)}-1=16\cdot4^{2k}-1=16(4^{2k}-1)+16-1=16(4^{2k}-1)+15$. Since $4^{2k}-1$ is divisible by $15$ (hypothesis) and $15$ is divisible by $15$, the sum is divisible by $15$.}
\inv{Why does $4^{2k}=16^k\equiv1\pmod{15}$ make this pattern unsurprising?}

%── B5 ──────────────────────────────────────────────────────────────────────────
\item Prove by induction that $5^n-1$ is divisible by $4$ for all positive integers $n$.

\ans{Proof by induction}
\method{Base case $n=1$: $5^1-1=4$, divisible by $4$.\\
Inductive step: assume $5^k-1=4m$ for some integer $m$. Then $5^{k+1}-1=5\cdot5^k-1=5(4m+1)-1=20m+4=4(5m+1)$, divisible by $4$. By induction, true for all positive integers $n$.}
\inv{Generalise: for which integer $a$ is $a^n-1$ divisible by $a-1$ for every positive integer $n$?}

%── B6 ──────────────────────────────────────────────────────────────────────────
\item Assuming $k!>2^k$ for some $k\geq4$, prove $(k+1)!>2^{k+1}$.

\ans{Proof}
\method{$(k+1)!=(k+1)\cdot k!>(k+1)\cdot2^k$ (using the hypothesis). Since $k\geq4$, $k+1\geq5>2$, so $(k+1)\cdot2^k>2\cdot2^k=2^{k+1}$. Combining, $(k+1)!>2^{k+1}$.}
\inv{Check the base case $k=4$: is $4!>2^4$?}

%── B7 ──────────────────────────────────────────────────────────────────────────
\item Prove by induction: $2^n\geq n+1$ for all non-negative integers $n$.

\ans{Proof by induction}
\method{Base case $n=0$: $2^0=1\geq0+1=1$. Holds (with equality).\\
Inductive step: assume $2^k\geq k+1$. Then $2^{k+1}=2\cdot2^k\geq2(k+1)=2k+2=(k+2)+k\geq k+2$ (since $k\geq0$). So $2^{k+1}\geq(k+1)+1$. By induction, true for all $n\geq0$.}
\inv{For which $n$ does equality $2^n=n+1$ actually hold?}

%── B8 ──────────────────────────────────────────────────────────────────────────
\item Assuming $\sum_{i=1}^k \frac1{i(i+1)}=\frac{k}{k+1}$, show that $\sum_{i=1}^{k+1}\frac1{i(i+1)}=\frac{k+1}{k+2}$.

\ans{$\frac{k}{k+1}+\frac1{(k+1)(k+2)}=\frac{k(k+2)+1}{(k+1)(k+2)}=\frac{(k+1)^2}{(k+1)(k+2)}=\frac{k+1}{k+2}$}
\method{$k(k+2)+1=k^2+2k+1=(k+1)^2$, which cancels one factor of $(k+1)$ from the denominator.}
\inv{This is a telescoping sum: write $\frac1{i(i+1)}=\frac1i-\frac1{i+1}$ and confirm the telescoping gives the same closed form directly.}

%── B9 ──────────────────────────────────────────────────────────────────────────
\item Prove by induction that the sum of the interior angles of a convex $n$-gon is $(n-2)\times180^\circ$ for all $n\geq3$.

\ans{Proof by induction}
\method{Base case $n=3$: a triangle has angle sum $180^\circ=(3-2)\times180^\circ$. Holds.\\
Inductive step: assume a convex $k$-gon has angle sum $(k-2)\times180^\circ$. A convex $(k+1)$-gon can be split, by drawing one diagonal from a vertex, into a convex $k$-gon and a triangle, whose angle sums combine to give the full $(k+1)$-gon's angle sum: $(k-2)\times180^\circ+180^\circ=(k-1)\times180^\circ=((k+1)-2)\times180^\circ$. By induction, true for all $n\geq3$.}
\inv{Why does this diagonal-splitting argument require the $(k+1)$-gon to be convex?}

%── B10 ──────────────────────────────────────────────────────────────────────────
\item A student wants to prove $P(n)$ for all $n\geq1$ using base cases $P(1)$ and $P(2)$, with inductive step ``$P(k)$ and $P(k+1)$ together imply $P(k+2)$'' for $k\geq1$. Explain why establishing \emph{both} $P(1)$ and $P(2)$ (not just $P(1)$) is necessary here.

\ans{The inductive step requires two consecutive prior cases to fire; without $P(2)$ independently established, the chain can never produce $P(2)$ or anything depending on it.}
\method{The step $P(k)\land P(k+1)\implies P(k+2)$ can only be applied once both $P(k)$ and $P(k+1)$ are already known. Starting from $P(1)$ alone gives no way to reach $P(2)$ --- the recursion needs two ``seeds'', exactly analogous to how the Fibonacci recurrence needs two initial terms, not one.}
\inv{How many base cases would be needed for a step of the form ``$P(k),P(k+1),P(k+2)\implies P(k+3)$''?}

\end{enumerate}

\section*{Section C}

\begin{enumerate}

%── C1 ──────────────────────────────────────────────────────────────────────────
\item A flawed ``proof'' that everyone in any group shares a favourite integer.

\ans{B}
\method{The overlap argument genuinely requires the two subgroups of size $k$ (formed by removing different people from a $(k+1)$-person group) to share at least one common member, which needs $k\geq2$. The transition from $P(1)$ to $P(2)$ uses $k=1$, where the ``two subgroups'' are single individuals that share no member --- so this specific step is invalid, and the whole chain never gets off the ground, regardless of how valid later steps might look in isolation.}
\inv{Confirm that the step from $P(2)$ to $P(3)$ (i.e.\ $k=2$) would genuinely be valid in isolation, if only it could be reached.}

%── C2 ──────────────────────────────────────────────────────────────────────────
\item Evaluate an inductive step for $n^3-n$ divisible by $6$ that appeals to ``the general theory'' instead of doing the algebra.

\ans{B}
\method{A correct inductive step must actually derive $(k+1)^3-(k+1)$ algebraically from $k^3-k=6m$: $(k+1)^3-(k+1)=k^3+3k^2+3k+1-k-1=(k^3-k)+3k^2+3k=6m+3k(k+1)$. Since $k(k+1)$ is a product of consecutive integers, it is even, so $3k(k+1)$ is divisible by $6$, and the whole expression is $6m+6(\text{integer})$, divisible by $6$. The original ``proof'' skips this derivation entirely, merely asserting the conclusion.}
\inv{Complete the missing step: why is $k(k+1)$ always even?}

%── C3 ──────────────────────────────────────────────────────────────────────────
\item Evaluate a base-case dispute for $n^2>3n+4$.

\ans{B}
\method{$n^2-3n-4=(n-4)(n+1)$, positive exactly when $n>4$ (for positive $n$). So $P(4)$: $16>16$ is false (equality, not strict), but $P(5)$: $25>19$ is true. The inductive step holds for all $k\geq5$: assuming $k^2-3k-4>0$ (i.e.\ $P(k)$), we have $(k+1)^2-3(k+1)-4=k^2+2k+1-3k-3-4=(k^2-3k-4)+(2k-2)$, a sum of a positive quantity (the inductive hypothesis) and another positive quantity ($2k-2\geq8>0$ for $k\geq5$), hence positive, giving $P(k+1)$. A false $P(n_0)$ at one specific point does not mean induction is impossible --- it only means that particular $n_0$ cannot serve as the base case.}
\inv{Verify directly that $P(5)\implies P(6)$: check $36>22$.}

%── C4 ──────────────────────────────────────────────────────────────────────────
\item Evaluate a weak-induction attempt to prove every $n\geq2$ is a product of primes.

\ans{B}
\method{The fatal gap: weak induction's hypothesis only concerns $k$, but factoring $k+1$ (when composite) requires information about its proper factors $a,b<k+1$, which are generally unrelated to $k$ itself. Strong induction fixes this by assuming the claim for every integer from $2$ to $k$, so that whichever factors $a,b$ of $k+1$ arise, both are already covered by the (stronger) inductive hypothesis.}
\inv{Write out the corrected strong-induction version of this proof in full.}

%── C5 ──────────────────────────────────────────────────────────────────────────
\item Evaluate whether a failed $P(1)$ means no base case exists for $\frac1{n+1}+\cdots+\frac1{2n}>\frac12$.

\ans{B}
\method{$P(1)$ gives exactly $\frac12$, so $\frac12>\frac12$ is false --- $n_0=1$ doesn't work. But $P(2)$: $\frac13+\frac14=\frac7{12}>\frac12$, true. In fact this sum is increasing in $n$ (each step replaces the smallest term with two smaller-but-larger-summing terms), so $n_0=2$ works as the base case.}
\inv{Show algebraically that the sum strictly increases when going from $n$ to $n+1$.}

%── C6 ──────────────────────────────────────────────────────────────────────────
\item Evaluate the ``$3$p/$5$p postage'' strong-induction step.

\ans{A}
\method{Since every increment of $1$ in the target amount can be achieved by adding one $3$p stamp to a construction that is $3$ less, the inductive step must reach back to $P(k-2)$, not $P(k)$ --- this is why three consecutive base cases ($8,9,10$) are needed to seed every residue class mod $3$ before the ``step back by $3$'' induction can run indefinitely.}
\inv{Which two-stamp combinations show that $11,12,13$ pence are also directly constructible, confirming the pattern from the base cases?}

%── C7 ──────────────────────────────────────────────────────────────────────────
\item Evaluate a finite-case ``proof'' of the alternating Fibonacci sum identity.

\ans{B}
\method{Direct check: $n=1$: $F_1=1=F_2$. $n=2$: $F_1+F_3=1+2=3=F_4$. $n=3$: $+F_5=3+5=8=F_6$. $n=4$: $+F_7=8+13=21=F_8$. All hold, but this is still only $4$ cases. The correct inductive step: assuming $F_1+\cdots+F_{2k-1}=F_{2k}$, add $F_{2k+1}$ to both sides: $F_1+\cdots+F_{2k-1}+F_{2k+1}=F_{2k}+F_{2k+1}=F_{2k+2}$ (by the Fibonacci recurrence itself), which is exactly the $n=k+1$ case.}
\inv{Why does this identity's inductive step require essentially no extra algebra, unlike most sum identities?}

%── C8 ──────────────────────────────────────────────────────────────────────────
\item Which statement precisely captures the inductive step?

\ans{C}
\method{The inductive step is itself a universally-quantified conditional statement: for every $k\geq1$, IF $P(k)$ holds THEN $P(k+1)$ holds. Options A, B, D, and E each misstate either the quantifier (existential instead of universal, or missing entirely) or the logical connective.}
\inv{Which option (D) describes something logically much weaker than the actual inductive step, and why is it too weak to be useful?}

\end{enumerate}

\section*{Section D}

\begin{enumerate}

%── D1 ──────────────────────────────────────────────────────────────────────────
\item $t=x+\frac1x$ integer $>2$; define $t_n=x^n+\frac1{x^n}$. \textit{\small(after BMO1 2015 Q4)}

\ans{(a) $t_{n+1}=t\cdot t_n-t_{n-1}$. (b) Proof by strong induction below.}
\method{(a) $t\cdot t_n=\left(x+\frac1x\right)\left(x^n+\frac1{x^n}\right)=x^{n+1}+x^{1-n}+x^{n-1}+\frac1{x^{n+1}}=\left(x^{n+1}+\frac1{x^{n+1}}\right)+\left(x^{n-1}+\frac1{x^{n-1}}\right)=t_{n+1}+t_{n-1}$. Rearranging gives $t_{n+1}=t\cdot t_n-t_{n-1}$.\\
(b) Base cases: $t_0=2$ (integer) and $t_1=t$ (integer, given). Strong inductive step: assume $t_{n-1}$ and $t_n$ are both integers for some $n\geq1$. Since $t$ is an integer and $t_{n+1}=t\cdot t_n-t_{n-1}$ is a sum/product of integers, $t_{n+1}$ is also an integer. By strong induction, $t_n$ is an integer for every positive integer $n$.}
\inv{Why does this proof genuinely need strong (two-step-back) induction rather than ordinary single-step induction?}

%── D2 ──────────────────────────────────────────────────────────────────────────
\item Prove every positive integer is a sum of distinct charming integers ($2$ or $3^i5^j$). \textit{\small(after BMO1 2016 Q6)}

\ans{Proof by strong induction below, via the key gap lemma.}
\method{\emph{Key lemma}: if $c<c'$ are consecutive charming integers (no charming integer strictly between them), then $c'<2c$. Spot-checking the sorted charming integers $1,2,3,5,9,15,25,27,45,75,81,\dots$ confirms every consecutive ratio is below $2$ (e.g.\ $9/5=1.8$, $15/9\approx1.67$, $27/25=1.08$); a full general proof requires bounding how far apart consecutive values of $3^i5^j$ can be, using that $5/3<2$ and $3<2\times2$ constrain how much any single-exponent increment can jump --- this gap-bound is the genuine crux of the problem's difficulty, exactly why it carries full marks as a BMO1 question rather than being a one-line observation.\\
Base case $n=1$: $1=3^05^0$ is itself charming, so it trivially ``sums'' to itself.\\
Strong inductive step: assume every positive integer less than $n$ (for some $n\geq2$) can be written as a sum of distinct charming integers. If $n$ is itself charming, it trivially works. Otherwise, let $c$ be the largest charming integer with $c<n$, and let $c'$ be the next charming integer above $c$. Since $c$ was chosen as the largest charming integer below $n$, we have $c'\geq n$; combined with the key lemma ($c'<2c$), this gives $n\leq c'<2c$, i.e.\ $n-c<c$. So $0<n-c<c<n$, and by the strong inductive hypothesis, $n-c$ is a sum of distinct charming integers, each of which is $\leq n-c<c$ and therefore cannot equal $c$ itself. Adding $c$ to this sum gives $n$ as a sum of distinct charming integers, completing the induction.}
\inv{Verify the base cases $n=1,2,3,4$ explicitly by hand, and check the gap lemma holds for the first ten charming integers listed above.}

%── D3 ──────────────────────────────────────────────────────────────────────────
\item Prove Bernoulli's Inequality: $(1+x)^n\geq1+nx$ for $x>-1$, $n$ a positive integer.

\ans{Proof by induction below; the condition $x>-1$ is required when multiplying the inequality by $(1+x)$.}
\method{Base case $n=1$: $(1+x)^1=1+x=1+1\cdot x$. Holds with equality.\\
Inductive step: assume $(1+x)^k\geq1+kx$ for some $k\geq1$. Since $x>-1$, we have $1+x>0$, so multiplying both sides of the hypothesis by $(1+x)$ \emph{preserves} the inequality direction: $(1+x)^{k+1}\geq(1+kx)(1+x)=1+x+kx+kx^2=1+(k+1)x+kx^2$. Since $kx^2\geq0$, this gives $(1+x)^{k+1}\geq1+(k+1)x+kx^2\geq1+(k+1)x$. By induction, true for all positive integers $n$.}
\inv{What would go wrong with the argument if $x<-1$ were allowed (i.e.\ $1+x<0$)?}

%── D4 ──────────────────────────────────────────────────────────────────────────
\item $a_1=2,a_{n+1}=a_n^2-a_n+1$. Prove $a_1,\dots,a_n$ pairwise coprime.

\ans{Proof by induction below (via the lemma $a_n-1=\prod_{i=1}^{n-1}a_i$).}
\method{\emph{Lemma}: $a_n-1=\prod_{i=1}^{n-1}a_i$ for all $n\geq2$, proved by induction on $n$.\\
Base case $n=2$: $a_2-1=(2^2-2+1)-1=3-1=2=a_1$. Holds.\\
Inductive step: assume $a_k-1=\prod_{i=1}^{k-1}a_i$. From the recurrence, $a_{k+1}-1=a_k^2-a_k=a_k(a_k-1)=a_k\prod_{i=1}^{k-1}a_i=\prod_{i=1}^{k}a_i$. Lemma holds for all $n\geq2$ by induction.\\
\emph{Pairwise coprimality}: take any $i<j\leq n$. By the lemma, $a_j-1=\prod_{m=1}^{j-1}a_m$, a product that includes the factor $a_i$ (since $i\leq j-1$). So $a_i\mid(a_j-1)$, i.e.\ $a_j\equiv1\pmod{a_i}$. Hence $\gcd(a_i,a_j)=\gcd(a_i,a_j-1+1)=\gcd(a_i,1)=1$ (using that $a_j-1$ is a multiple of $a_i$). Since $i,j$ were arbitrary, $a_1,\dots,a_n$ are pairwise coprime.}
\inv{Compute $a_1,a_2,a_3,a_4$ explicitly and verify they are pairwise coprime by direct inspection.}

%── D5 ──────────────────────────────────────────────────────────────────────────
\item Prove $\sum_{k=1}^n\frac1{k^2}<2-\frac1n$ for all integers $n\geq2$.

\ans{Proof by induction below.}
\method{Base case $n=2$: LHS $=1+\frac14=\frac54$. RHS $=2-\frac12=\frac32$. Since $\frac54<\frac32$, holds.\\
Inductive step: assume $\sum_{k=1}^n\frac1{k^2}<2-\frac1n$ for some $n\geq2$. Then $\sum_{k=1}^{n+1}\frac1{k^2}<2-\frac1n+\frac1{(n+1)^2}$. It suffices to show $-\frac1n+\frac1{(n+1)^2}\leq-\frac1{n+1}$, i.e.\ $\frac1{(n+1)^2}\leq\frac1n-\frac1{n+1}=\frac1{n(n+1)}$. This holds because $(n+1)^2>n(n+1)$ (as $n+1>n$), so $\frac1{(n+1)^2}<\frac1{n(n+1)}$, strictly. Combining, $\sum_{k=1}^{n+1}\frac1{k^2}<2-\frac1n+\frac1{(n+1)^2}<2-\frac1n+\frac1{n(n+1)}=2-\frac1{n+1}$. By induction, true for all $n\geq2$.}
\inv{What does this inequality suggest about the value of $\sum_{k=1}^\infty\frac1{k^2}$ as $n\to\infty$?}

\end{enumerate}

\end{document}
